\begin{notechapter}
\chapter{Bayes Decision Theory}\label{ch:bayes-decision-theory}
\chaptersubtitle{The best possible predictor}

\begin{chapterintro}
Suppose for a moment that the distribution $P$ of $(X,Y)$ were known. Which
predictor should we use? This chapter answers the question for every common
loss, and in doing so defines the benchmark against which every learning
algorithm is measured. You will learn
\begin{itemize}[nosep]
  \item the \emph{pointwise principle}: the best predictor minimizes the
        expected loss separately at each input;
  \item that under the 0--1 loss the best classifier picks the most probable
        class (the \emph{Bayes classifier}), and how it relates to the MAP and
        maximum-likelihood rules;
  \item how unequal error costs move the decision threshold;
  \item that under the squared loss the best predictor is the conditional
        mean $\E[Y \mid X]$, with two proofs, and what happens when the model
        class is too small to contain it;
  \item the best predictors for the absolute loss and the log-loss.
\end{itemize}
\prerequisites{\cref{ch:learning-problem}; Bayes' rule; conditional
expectation and the tower property $\E[Z] = \E\bigl[\E[Z\mid X]\bigr]$.}
\end{chapterintro}

\section{The pointwise principle}\label{bayes-decision-theory:sec:pointwise}

Recall that the risk of a predictor is $\Risk(f) = \E[\ell(f(X),Y)]$. We look
for a predictor of smallest risk among \emph{all} functions $\Xs \to \Ys$.

\begin{definition}[Bayes predictor and Bayes risk]
The \emph{Bayes risk} is the smallest risk achievable by any predictor,
\[
  R^* = \inf_{f} \Risk(f),
\]
the infimum being over all (measurable) functions $f:\Xs\to\Ys$. A predictor
$f^*$ with $\Risk(f^*) = R^*$ is a \emph{Bayes predictor} (\emph{Bayes
classifier} in classification).
\end{definition}

The key observation is that the choice of $f(x)$ at one input does not
constrain the choice at another, so the global problem splits into one small
problem per input.

\begin{theorem}[Pointwise minimization]\label{bayes-decision-theory:thm:pointwise}
For $x \in \Xs$ and a candidate prediction $a \in \Ys$, let
\[
  r(a \mid x) = \E\bigl[\ell(a, Y) \mid X = x\bigr]
\]
be the \emph{conditional risk} of predicting $a$ at $x$. If $f^*(x)$ minimizes
$a \mapsto r(a\mid x)$ for every $x$, then $f^*$ is a Bayes predictor, and
\[
  R^* = \E\Bigl[\min_{a} r(a \mid X)\Bigr].
\]
\end{theorem}

\begin{proof}
For any predictor $f$, the tower property gives
\[
  \Risk(f) = \E\Bigl[\E\bigl[\ell(f(X), Y)\mid X\bigr]\Bigr] = \E\bigl[r(f(X)\mid X)\bigr].
\]
By the choice of $f^*$, $r(f(x)\mid x) \geq r(f^*(x)\mid x)$ for every $x$.
Taking expectations over $X$ preserves the inequality, so
$\Risk(f) \geq \Risk(f^*)$ for every $f$.
\end{proof}

\begin{keyidea}
To find the best predictor, fix the input $x$, look at the conditional
distribution of $Y$ given $X = x$, and choose the single value $a$ that has
the smallest expected loss under that distribution. Only $P(Y\mid X)$ matters;
the marginal $P(X)$ only weights the inputs when computing $R^*$.
\end{keyidea}

The rest of the chapter applies this principle to specific losses.

\section{Classification with the 0--1 loss}\label{bayes-decision-theory:sec:01}

Let $\Ys = \{1, \dots, K\}$ and write
$\eta_k(x) = P(Y = k \mid X = x)$ for the posterior class probabilities.
Under the 0--1 loss the conditional risk of predicting class $a$ is the
probability that the true class is something else:
\[
  r(a \mid x) = \sum_{k} \bbone[a \neq k]\, \eta_k(x) = 1 - \eta_a(x).
\]
Minimizing it means maximizing $\eta_a(x)$.

\begin{theorem}[Bayes classifier]\label{bayes-decision-theory:thm:bayes-classifier}
Under the 0--1 loss, the classifier
\[
  f^*(x) = \argmax_{k \in \Ys} P(Y = k \mid X = x)
\]
(ties broken arbitrarily) is a Bayes classifier, and the Bayes risk is
\[
  R^* = \E\Bigl[1 - \max_k \eta_k(X)\Bigr].
\]
In the binary case, with $\eta(x) = P(Y=1\mid X=x)$, the Bayes classifier
predicts $1$ exactly when $\eta(x) > \frac12$, and $R^* = \E\bigl[\min(\eta(X), 1-\eta(X))\bigr]$.
\end{theorem}

\begin{proof}
Apply \cref{bayes-decision-theory:thm:pointwise} with
$r(a\mid x) = 1 - \eta_a(x)$. In the binary case
$\max(\eta, 1-\eta) > \frac12$ is attained by class 1 iff $\eta > \frac12$,
and $1 - \max(\eta, 1-\eta) = \min(\eta, 1-\eta)$.
\end{proof}

The Bayes risk is zero only if $Y$ is a deterministic function of $X$. In
general it is positive: it is the \emph{irreducible error} caused by the
information missing from $X$.

\begin{example}[The die, continued]\label{bayes-decision-theory:ex:die}
In \cref{learning-problem:ex:die}, the posteriors are
\[
  P(Y=1 \mid X=0) = \frac{2/6}{3/6} = \frac23, \qquad
  P(Y=1 \mid X=1) = \frac{1/6}{3/6} = \frac13 .
\]
The Bayes classifier predicts $f^*(0) = 1$ and $f^*(1) = 0$ --- it is the
classifier ``flip $X$'' found by enumeration in \cref{ch:learning-problem}.
At each input the conditional error is $1 - \frac23 = \frac13$, so
$R^* = \frac12 \cdot \frac13 + \frac12 \cdot \frac13 = \frac13$. One throw in
three is misclassified whatever we do, because parity does not determine
primality.
\end{example}

How much worse than optimal is another classifier? In the binary case there is
an exact formula.

\begin{proposition}[Excess risk of a binary classifier]\label{bayes-decision-theory:prop:excess}
For any classifier $g : \Xs \to \{0,1\}$,
\[
  \Risk(g) - R^* = \E\Bigl[\,\abs{2\eta(X) - 1}\;\bbone\bigl[g(X) \neq f^*(X)\bigr]\Bigr] \;\geq\; 0 .
\]
Consequently $\Risk(g) = R^*$ if and only if $g$ agrees with $f^*$ for almost
every input at which $\eta(X) \neq \frac12$.
\end{proposition}

\begin{proof}
At input $x$, the conditional error of $g$ is $\eta(x)$ if $g(x) = 0$ and
$1 - \eta(x)$ if $g(x) = 1$. Where $g(x) = f^*(x)$ the two conditional errors
coincide. Where $g(x) \neq f^*(x)$, $g$ picks the less likely class and its
conditional error exceeds that of $f^*$ by
$\max(\eta, 1 - \eta) - \min(\eta, 1-\eta) = \abs{2\eta(x) - 1}$. Taking the
expectation over $X$ gives the formula; the integrand is non-negative and
vanishes exactly where $g = f^*$ or $\eta = \frac12$.
\end{proof}

Disagreeing with the Bayes classifier is cheap where the posterior is close to
$\frac12$ (both classes are almost equally likely) and expensive where it is
close to $0$ or $1$.

\section{Bayes, MAP and maximum-likelihood classifiers}\label{bayes-decision-theory:sec:map-ml}

The posterior is rarely modelled directly in generative approaches; one
models instead the class-conditional distributions $P(X = x \mid Y = k)$ (the
\emph{likelihoods}) and the class frequencies $P(Y = k)$ (the \emph{priors}).
Bayes' rule connects them:
\[
  P(Y = k \mid X = x) = \frac{P(X = x \mid Y = k)\, P(Y = k)}{P(X = x)} .
\]
The denominator does not depend on $k$, so it plays no role in the $\argmax$.
This gives two classical decision rules.

\begin{definition}[MAP and ML classifiers]
The \emph{maximum a posteriori} (MAP) classifier and the \emph{maximum-likelihood}
(ML) classifier are
\[
  f_{\mathrm{MAP}}(x) = \argmax_k\, P(X = x \mid Y = k)\,P(Y = k), \qquad
  f_{\mathrm{ML}}(x) = \argmax_k\, P(X = x \mid Y = k).
\]
\end{definition}

The relations between the three rules are now immediate.
\begin{itemize}
  \item \textbf{MAP = Bayes under the 0--1 loss.} $f_{\mathrm{MAP}}$ maximizes
        the posterior, which is exactly the Bayes classifier of
        \cref{bayes-decision-theory:thm:bayes-classifier}. ``Minimizing the
        probability of error'' and ``maximizing the posterior'' are the same
        objective; they only part ways when the costs of errors differ
        (\cref{bayes-decision-theory:sec:costs}).
  \item \textbf{ML = MAP with uniform priors.} If $P(Y=k) = 1/K$ for all $k$,
        the prior factor is constant and drops out of the $\argmax$. Otherwise
        the ML rule ignores how common each class is, and is in general
        suboptimal.
\end{itemize}

\begin{example}[When the prior changes the decision]
Two classes have priors $P(Y=0) = 0.7$ and $P(Y=1) = 0.3$. We observe an input
$x$ and consider two possible likelihood models.

\emph{Case 1:} $P(x \mid Y=0) = 0.2$, $P(x\mid Y=1) = 0.8$. The unnormalized
posteriors are $0.2 \times 0.7 = 0.14$ and $0.8 \times 0.3 = 0.24$; after
normalization $P(Y = 1\mid x) = 0.24/0.38 \approx 0.63$. Both MAP and ML
predict class 1.

\emph{Case 2:} $P(x \mid Y=0) = 0.2$, $P(x\mid Y=1) = 0.4$. The unnormalized
posteriors are $0.14$ and $0.12$, so $P(Y=1\mid x) = 0.12/0.26 \approx 0.46$.
MAP predicts class 0; ML still predicts class 1, because $x$ is twice as
likely under class 1. Here the evidence is not strong enough to overturn the
prior, and the ML rule makes the wrong call more than half of the time.
\end{example}

\section{Unequal costs of errors}\label{bayes-decision-theory:sec:costs}

Declaring a sick patient healthy is worse than the reverse. Such asymmetries
are encoded in a \emph{cost matrix} $C$, where $C_{k a}$ is the cost of
predicting class $a$ when the true class is $k$; usually $C_{kk} = 0$. The loss
is $\ell(a, k) = C_{k a}$ and the risk becomes
\[
  \Risk(f) = \E\bigl[C_{Y, f(X)}\bigr].
\]
The pointwise principle gives the optimal rule at once.

\begin{theorem}[Cost-sensitive Bayes classifier]\label{bayes-decision-theory:thm:costs}
The classifier
\[
  f^*(x) = \argmin_{a \in \Ys} \sum_{k\in\Ys} \eta_k(x)\, C_{k a}
\]
minimizes $\E[C_{Y,f(X)}]$. In the binary case with $C_{00} = C_{11} = 0$, it
predicts class $1$ exactly when
\[
  \eta(x) > \frac{C_{01}}{C_{01} + C_{10}} .
\]
\end{theorem}

\begin{proof}
The conditional risk of predicting $a$ is $r(a\mid x) = \sum_k \eta_k(x) C_{ka}$;
apply \cref{bayes-decision-theory:thm:pointwise}. In the binary case,
predicting $0$ costs $C_{10}$ when $Y = 1$, so $r(0\mid x) = C_{10}\,\eta(x)$;
predicting $1$ costs $C_{01}$ when $Y = 0$, so $r(1\mid x) = C_{01}(1 - \eta(x))$.
Class 1 is preferred when $C_{01}(1-\eta) < C_{10}\,\eta$, that is when
$\eta > C_{01}/(C_{01}+C_{10})$.
\end{proof}

With equal costs the threshold is $\frac12$ and we recover the Bayes
classifier. Raising the cost of missing class 1 lowers the threshold, so
class 1 is predicted more readily.

\begin{example}[Medical screening]\label{bayes-decision-theory:ex:medical}
Let class 1 be ``sick''. A false alarm costs $C_{01} = 1$ (an extra test); a
missed case costs $C_{10} = 10$. The threshold becomes
$1/(1 + 10) = 1/11 \approx 0.09$: the patient is flagged as soon as the
probability of disease exceeds 9\%, not 50\%.
\Cref{bayes-decision-theory:fig:costs} shows the two expected costs as
functions of $\eta$; the optimal decision is the lower line.
\end{example}

\begin{figure}[ht]
\centering
\begin{tikzpicture}
\begin{axis}[width=0.5\textwidth, height=5.2cm, xmin=0, xmax=1, ymin=0, ymax=1.6,
  xlabel={$\eta(x) = P(Y=1\mid X=x)$}, ylabel={expected cost},
  xtick={0, 0.0909, 0.5, 1}, xticklabels={$0$, $\frac{1}{11}$, $\frac12$, $1$},
  legend style={at={(1.03,0.5)}, anchor=west, font=\small}, axis lines=left]
  \addplot[thick, bookaccent, domain=0:0.16] {10*x};
  \addlegendentry{predict 0: $10\,\eta$}
  \addplot[thick, bookpitfall, domain=0:1] {1-x};
  \addlegendentry{predict 1: $1-\eta$}
  \addplot[dashed, gray, domain=0:1] {x};
  \addlegendentry{predict 0, equal costs: $\eta$}
  \draw[dotted] (axis cs:0.0909,0) -- (axis cs:0.0909,0.909);
  \draw[dotted] (axis cs:0.5,0) -- (axis cs:0.5,0.5);
\end{axis}
\end{tikzpicture}
\caption{Expected cost of each decision in \cref{bayes-decision-theory:ex:medical}.
With equal costs the lines $\eta$ and $1-\eta$ cross at $\frac12$; making a
missed case ten times more costly moves the crossing point to $\frac1{11}$.}
\label{bayes-decision-theory:fig:costs}
\end{figure}

In practice, costs are taken into account in one of four ways:
\begin{itemize}
  \item \textbf{threshold adjustment}: train a probabilistic classifier as
        usual and threshold its estimate of $\eta$ at
        $C_{01}/(C_{01}+C_{10})$ instead of $\frac12$ --- the direct
        application of \cref{bayes-decision-theory:thm:costs};
  \item \textbf{weighted losses}: multiply the training loss of each example
        by the cost of misclassifying its class;
  \item \textbf{resampling}: over-sample the costly class (or under-sample the
        other), which has a similar effect on the learned model;
  \item \textbf{cost-sensitive algorithms}, such as decision trees or SVMs
        whose splitting criterion or margin penalty is weighted by $C$.
\end{itemize}

\section{Regression with the squared loss}\label{bayes-decision-theory:sec:squared}

Let $\Ys = \R$ and $\ell(\hat y, y) = (\hat y - y)^2$. Write
$m(x) = \E[Y \mid X = x]$ for the \emph{regression function}.

\begin{theorem}[The conditional mean is optimal]\label{bayes-decision-theory:thm:mse}
Assume $\E[Y^2] < \infty$. For every predictor $f$ with $\E[f(X)^2]<\infty$,
\begin{equation}\label{bayes-decision-theory:eq:pythagoras}
  \E\bigl[(f(X) - Y)^2\bigr]
  = \underbrace{\E\bigl[(Y - m(X))^2\bigr]}_{\text{does not depend on } f}
  + \E\bigl[(f(X) - m(X))^2\bigr].
\end{equation}
Hence $f^* = m$ is a Bayes predictor, the Bayes risk is the expected
conditional variance $R^* = \E[\Var(Y\mid X)]$, and any other Bayes predictor
equals $m$ almost surely.
\end{theorem}

\emph{First proof (calculus).} By the pointwise principle we minimize, for a
fixed $x$, the conditional risk
\[
  r(a\mid x) = \E[(a - Y)^2 \mid X = x] = a^2 - 2a\,m(x) + \E[Y^2\mid X=x],
\]
a convex parabola in $a$. Its derivative $2a - 2m(x)$ vanishes at
$a = m(x)$, which is therefore the unique minimizer.
\hfill$\square$

\emph{Second proof (orthogonal decomposition).} Write
$f(X) - Y = \bigl(f(X) - m(X)\bigr) + \bigl(m(X) - Y\bigr)$ and expand the
square. The cross term vanishes:
\[
  \E\bigl[(f(X) - m(X))(m(X) - Y)\bigr]
  = \E\Bigl[(f(X) - m(X))\,\underbrace{\E[m(X) - Y \mid X]}_{=\,0}\Bigr] = 0,
\]
because $f(X) - m(X)$ is a function of $X$ and can be taken out of the
conditional expectation. This proves \eqref{bayes-decision-theory:eq:pythagoras}.
The last term is non-negative and is zero if and only if $f(X) = m(X)$ almost
surely. Finally $\E[(Y-m(X))^2] = \E\bigl[\E[(Y-m(X))^2\mid X]\bigr] =
\E[\Var(Y\mid X)]$.
\hfill$\square$

\begin{intuition}[Pythagoras in a space of random variables]
Equation \eqref{bayes-decision-theory:eq:pythagoras} is Pythagoras' theorem.
Random variables with finite second moment form a vector space with inner
product $\inner{U}{V} = \E[UV]$. The conditional mean $m(X)$ is the orthogonal
projection of $Y$ onto the subspace of functions of $X$, the residual
$Y - m(X)$ is orthogonal to that subspace, and $f(X) - m(X)$ lies inside it.
The squared distance from $Y$ to any $f(X)$ is the squared distance to the
projection plus the squared distance from the projection to $f(X)$.
\end{intuition}

The second proof gives more than optimality: it says that the excess risk of
$f$ is exactly its mean squared distance to $m$. It also carries over verbatim
to vector outputs, with $\norm{f(X) - Y}^2$ and inner products in place of
products. In \emph{Deep Learning} by Goodfellow, Bengio and Courville
(Section~6.2.1.2) the same result is derived ``by calculus of variations'';
the first proof above is that argument made elementary.

\begin{pitfall}[``Unique'' means unique almost everywhere]
A Bayes predictor can be changed arbitrarily on a set of inputs of probability
zero without changing its risk. Statements such as ``$m$ is the unique
optimal predictor'' must be read as ``unique up to a set of probability
zero''.
\end{pitfall}

\subsection{When the model class is too small}\label{bayes-decision-theory:sec:projection}

Learning algorithms do not search over all functions but over a class $\Hs$,
for instance linear functions. Equation \eqref{bayes-decision-theory:eq:pythagoras}
holds for every $f$, so minimizing the risk over $\Hs$ is the same as
minimizing $\E[(f(X)-m(X))^2]$: \emph{the best predictor in $\Hs$ is the
element of $\Hs$ closest to $m$} in mean square. When $\Hs$ is a closed linear
space this is the orthogonal projection of $m$ onto $\Hs$.

\begin{example}[The best line through a parabola]\label{bayes-decision-theory:ex:parabola}
Let $X$ be uniform on $[-1,1]$ and $Y = X^2 + \varepsilon$, where
$\varepsilon$ is independent noise with mean $0$ and variance $s^2$. Then
$m(x) = x^2$ and $R^* = s^2$.

Among linear predictors $f(x) = a + bx$, the risk is minimized by
$b = \Cov(X, Y)/\Var(X)$ and $a = \E[Y] - b\,\E[X]$. Here
$\Cov(X,Y) = \E[X^3] = 0$ by symmetry, so $b = 0$ and
$a = \E[X^2] = \frac13$: the best line is \emph{horizontal}, although $Y$
depends strongly on $X$. Its risk is
\[
  s^2 + \E\bigl[(X^2 - \tfrac13)^2\bigr] = s^2 + \E[X^4] - \tfrac19
  = s^2 + \tfrac15 - \tfrac19 = s^2 + \tfrac{4}{45}.
\]
The extra $\frac4{45} \approx 0.089$ is the price of the restriction to
lines; \cref{ch:erm} calls it the \emph{approximation error}.
\end{example}

For linear regression on $\R^d$ the best linear predictor is unique as a
function whenever $\E[XX^\top]$ is invertible; with collinear features its
\emph{coefficients} are not unique, though the predictions are. For
non-convex classes such as neural networks, several different functions may
achieve the minimum.

\section{Other losses}\label{bayes-decision-theory:sec:other}

\paragraph{Absolute loss.} For $\ell(\hat y, y) = \abs{\hat y - y}$, the
conditional risk $a \mapsto \E[\abs{a - Y}\mid X=x]$ is minimized by any
\emph{median} of the conditional distribution of $Y$ (\cref{bayes-decision-theory:ex:median}).
Medians are insensitive to extreme values, which makes the absolute loss
robust to outliers.

\paragraph{Log-loss.} Let the prediction at $x$ be a probability vector
$q(x) = (q_1(x), \dots, q_K(x))$ and $\ell(q, y) = -\log q_y$.

\begin{proposition}[The log-loss is minimized by the true posterior]\label{bayes-decision-theory:prop:logloss}
For every $x$ and every probability vector $q$,
\[
  \E\bigl[-\log q_Y \mid X = x\bigr]
  = H\bigl(\eta(x)\bigr) + \mathrm{KL}\bigl(\eta(x)\,\Vert\, q\bigr)
  \;\geq\; H\bigl(\eta(x)\bigr),
\]
where $H(\eta) = -\sum_k \eta_k \log \eta_k$ is the entropy and
$\mathrm{KL}(\eta\Vert q) = \sum_k \eta_k \log(\eta_k/q_k)$ the Kullback--Leibler
divergence. Equality holds if and only if $q = \eta(x)$: the Bayes predictor
is the conditional distribution itself, $f^*(x) = \eta(x)$.
\end{proposition}

\begin{proof}
Adding and subtracting $\log \eta_k$,
$-\sum_k \eta_k \log q_k = -\sum_k \eta_k \log \eta_k + \sum_k \eta_k \log(\eta_k / q_k)$.
It remains to show $\mathrm{KL}(\eta\Vert q)\geq 0$ with equality iff $q=\eta$
(Gibbs' inequality). Since $\log$ is concave, Jensen's inequality gives
\[
  -\mathrm{KL}(\eta\Vert q) = \sum_{k:\eta_k>0} \eta_k \log\frac{q_k}{\eta_k}
  \leq \log \sum_{k:\eta_k>0} \eta_k \frac{q_k}{\eta_k}
  \leq \log 1 = 0,
\]
with equality iff $q_k/\eta_k$ is constant on $\{\eta_k > 0\}$ and
$q$ puts no mass elsewhere, that is iff $q = \eta$.
\end{proof}

Unlike the 0--1 loss, which only asks for the most likely class, the log-loss
rewards a predictor for getting the \emph{whole} posterior right. This is why
models trained with cross-entropy produce probabilities, which can then be
thresholded for any cost matrix by \cref{bayes-decision-theory:thm:costs}.
\Cref{bayes-decision-theory:tab:summary} collects the results of the chapter.

\begin{table}[ht]
\centering
\caption{Bayes predictors for common losses.}
\label{bayes-decision-theory:tab:summary}
\begin{tabular}{@{}lll@{}}
\toprule
Loss & Bayes predictor $f^*(x)$ & Bayes risk $R^*$ \\
\midrule
0--1 & $\argmax_k \eta_k(x)$ & $\E[1 - \max_k \eta_k(X)]$ \\
cost matrix $C$ & $\argmin_a \sum_k \eta_k(x) C_{ka}$ & $\E[\min_a \sum_k \eta_k(X) C_{ka}]$ \\
squared & $\E[Y\mid X=x]$ & $\E[\Var(Y\mid X)]$ \\
absolute & a median of $Y$ given $X = x$ & $\E[\,\abs{Y - \mathrm{med}(Y\mid X)}\,]$ \\
log-loss & $\eta(x) = P(Y = \cdot \mid X = x)$ & $\E[H(\eta(X))]$ \\
\bottomrule
\end{tabular}
\end{table}

\begin{chaptersummary}
\begin{itemize}[nosep]
  \item The Bayes predictor minimizes the conditional risk
        $\E[\ell(a, Y)\mid X=x]$ separately at each $x$; it depends only on
        $P(Y\mid X)$.
  \item Under the 0--1 loss it is the most probable class (the MAP rule); the
        ML rule coincides with it only when the classes are equally likely.
  \item The Bayes risk $R^*$ is the irreducible error; it is zero only when
        $Y$ is a function of $X$. For binary classification
        $\Risk(g) - R^* = \E[\abs{2\eta-1}\bbone[g\neq f^*]]$.
  \item Unequal costs move the binary threshold from $\frac12$ to
        $C_{01}/(C_{01}+C_{10})$.
  \item Under the squared loss the Bayes predictor is $\E[Y\mid X]$, and the
        excess risk of any $f$ is $\E[(f(X) - \E[Y\mid X])^2]$; restricting to a
        class $\Hs$ selects the element of $\Hs$ closest to $\E[Y\mid X]$.
  \item The absolute loss leads to the conditional median, the log-loss to the
        full conditional distribution.
\end{itemize}
\end{chaptersummary}

\section*{Exercises}
\addcontentsline{toc}{section}{Exercises}

\begin{exercise}\label{bayes-decision-theory:ex:median}
Let $Z$ be a real random variable with $\E\abs{Z} < \infty$ and $g(a) = \E\abs{a - Z}$.
Show that for $h > 0$,
$g(a+h) - g(a) = h\,\bigl(2\Prob(Z \leq a) - 1\bigr) + 2\,\E\bigl[(a + h - Z)\bbone[a < Z < a+h]\bigr]$,
and deduce that $g$ is minimized at any median of $Z$.
\hint{split according to $Z \leq a$, $a < Z < a+h$ and $Z \geq a+h$.}
\end{exercise}

\begin{exercise}
Two classes are equally likely and the class-conditional densities of
$X \in \R$ are $\mathcal N(-1, 1)$ and $\mathcal N(1,1)$. Show that the Bayes
classifier is $f^*(x) = \bbone[x > 0]$ and that $R^* = \Phi(-1) \approx 0.159$,
where $\Phi$ is the standard normal distribution function. How does the
threshold move if $P(Y=1) = 0.2$?
\end{exercise}

\begin{exercise}
Generalize \cref{bayes-decision-theory:thm:costs} to $C_{00}, C_{11}$ not
necessarily zero, assuming $C_{10} > C_{11}$ and $C_{01} > C_{00}$.
\end{exercise}

\begin{exercise}
In \cref{bayes-decision-theory:ex:parabola}, find the best \emph{quadratic}
predictor $a + bx + cx^2$ and its risk. What is its approximation error?
\end{exercise}
\end{notechapter}
